\chapter{Relational Dynamics of Frames: Autoinertia, Ambivalence, and Mach's Principle}
\label{chap:autoinercia-mach}

Classical inertia characterizes the persistence of a state of motion.
Mach's question \cite{MachMechanics} shifts the emphasis from the isolated body toward the
relation between that body and the material totality. In genealogical
language, the same question is formulated without separating the observer from the
boundary it reads: a local frame is the restriction of a globally compatible
state, and the inertial response is Machian when it depends on the totality through
a well-defined boundary trace.

The word \emph{self-inertial} does not designate a system disconnected from the rest.
It designates a section that coherently transports the frame induced by its
own insertion into the whole. Its persistence is internal to the global state
and is verified by means of a connection.

\section{Totality, Trace, and Local Response}

Let \(\Omega_{\rm tot}\) be the space of global sections and
\(b:\Omega_{\rm tot}\to\mathcal B\) its trace on a boundary. For a
region or vertex \(v\), let
\[
 \mathcal I_v:\Omega_{\rm tot}\longrightarrow\mathcal Y_v
\]
the local inertial response.

\begin{definition}[Machian Factorization]
The response \(\mathcal I_v\) is Machian with respect to \(b\) when there exists
an induced map
\(\overline{\mathcal I}_v:\operatorname{im}(b)\to\mathcal Y_v\) such that
\begin{equation}
 \mathcal I_v=\overline{\mathcal I}_v\circ b.
 \label{eq:factorizacion-machiana}
\end{equation}
\end{definition}

\begin{theorem}[Completeness Factorization Criterion]
\label{thm:criterio-mach}
There exists a unique map
\(\overline{\mathcal I}_v:\operatorname{im}(b)\to\mathcal Y_v\) that satisfies
\eqref{eq:factorizacion-machiana} if and only if \(\mathcal I_v\) is constant
over each fiber of \(b\):
\[
 b(x)=b(y)\quad\Longrightarrow\quad
 \mathcal I_v(x)=\mathcal I_v(y).
\]
If \(b\) is surjective, \(\operatorname{im}(b)=\mathcal B\) and the factorized
map is uniquely defined on all of \(\mathcal B\).
\end{theorem}

\begin{proof}
Necessity is immediate. For sufficiency, given \(\beta\in\operatorname{im}(b)\), we define \(\overline{\mathcal I}_v(\beta)=\mathcal I_v(x)\) for any \(x\in b^{-1}(\beta)\). Constancy on the fibers ensures that the definition is independent of the representative; the factorization equality itself forces its value at each \(\beta\in\operatorname{im}(b)\) and thereby ensures uniqueness. If \(b\) is surjective, this domain coincides with \(\mathcal B\).
\end{proof}

The criterion gives a positive formulation of Mach's principle: the global
boundary contains exactly the information of the totality that governs the
local response. The interior state may conserve more genealogy;
inertia depends on the class published by the trace.

\section{Self-inertial sections as geodesics of the joint state of system, observer, and memory: \(D_\gamma s_\gamma=0\)}

Let \(\mathcal E\to\Gamma\) to vector bundle of local states along
to route \(\gamma\), and is \(D\) to linear connection induced by
TPK transport.

\begin{definition}[Self-inertial Section]
A persistent section \(s_\gamma\) is self-inertial on a segment if
\begin{equation}
 D_\gamma s_\gamma=0.
 \label{eq:autoinercia-paralela}
\end{equation}
Equivalently and discretely, the state at each vertex is the parallel
transport of the preceding state by the route morphism.
\end{definition}

Equation \eqref{eq:autoinercia-paralela} preserves orientation, memory, and
signature in the sense determined by the connection. A coordinate acceleration
can coexist with parallel transport when the reference frame also
changes. The distinction between inertial motion and force is thus tied to
the connection, not to the appearance of the trajectory in an isolated chart.

\begin{proposition}[Persistence and Holonomy]
\label{prop:autoinercia-holonomia}
Let \(\gamma\) be a loop and \(\Hol_D(\gamma)\) its holonomy. A
self-inertial section returns exactly to its initial state if and only if it belongs to the
fixed subspace of \(\Hol_D(\gamma)\). If the visible projection returns and the
initial state lies outside
\(\operatorname{Fix}(\Hol_D(\gamma))\), the coordinate returns and the
state changes.
\end{proposition}

This result links self-inertia and monodromy. Observable periodicity does not
imply closure of the lift; the deficit or excess of holonomy records the
accumulated response of the frame.

\section{Ambivalence Principle: Bidirectional Involution of the Readers π/φ² and φ/π²}

The HMT cell admits two complementary readings: areal and volumetrica. Let
\(q_A\) the boundary or area reader and \(q_V\) the interior or
volume reader. The transport \(T_{AV}\) relates the two without identifying them.

\begin{principle}[Areal-volumetric Ambivalence]
A complete physical realization conserves the pair
\begin{equation}
 \bigl(q_A(x),q_V(x),T_{AV}(x)\bigr).
 \label{eq:ambivalencia-av}
\end{equation}
The areal reading expresses how the state impinges on a boundary; the
volumetric reading expresses how it occupies and transports the interior. The third datum
preserves the transformation between them.
\end{principle}

In Dimensional Mechanics, the areal--volumetric relation organizes the
HMT extension of the equivalence principle. A single section may be read
as a boundary response and as an interior density. When both readings
are transported by the same local connection, the inertial and
gravitational descriptions belong to a single state class.

\begin{criterion}[Typed local equivalence]
The identification of an inertial response with a gravitational response
is proved in a region \(U\) when:
\begin{enumerate}
 \item \(q_A\) and \(q_V\) are defined on the same section;
 \item there exists an intertwiner \(T_{AV}\) that preserves its units and
       orientations;
 \item the connection transporting both readers coincides on \(U\);
 \item the observational difference vanishes in the declared local frame.
\end{enumerate}
\end{criterion}

The formulation preserves the equivalence and shows which data realize it. Its
application to the electron belongs to the first physical closure of the mass
law; the present annex uses only the geometric scheme.

\section{Causal Coordination Between Inertial Transport, Areal–Volumetric Double Reading, and Machian Factorization}

The three principles are causally ordered:
\begin{equation}
 \boxed{
 \begin{aligned}
 &\text{persistence of the section}
   &&\longrightarrow&& D_\gamma s=0,\\
 &\text{double reading of the same state}
   &&\longrightarrow&& (q_A,q_V,T_{AV}),\\
 &\text{dependence of the local response on the whole}
   &&\longrightarrow&& \mathcal I_v=\overline{\mathcal I}_v\circ b.
 \end{aligned}}
 \label{eq:triptico-inercia-mach}
\end{equation}
The principle of inertia fixes the transport; Ambivalence fixes the two
physical readings; Mach fixes the global quotient that governs the local frame.
Together they produces to relational theory of motion in the that the observer,
the apparatus, and the body belong to to common section.

\section{Measurement of motion and frame coupling}

Measuring velocity or acceleration requires comparing the system with a clock and a
ruler. Those objects are part of the apparatus. Let \(F_S\) be the system frame
and \(F_M\) the observer frame. The measurement coupling constructs a
relative transformation
\[
 G_{MS}=F_M^{-1}F_S.
\]
The published velocity, acceleration, and phase are invariants or
components of \(G_{MS}\), not properties without reference. Repetition of the
experiment preserves the context by jointly transporting apparatus and
boundary.

This formulation recovers Weyl's intuition \cite{Weyl1927}: measuring an amplitude or a
length entails coupling a gauge rule. HMT adds the genealogy of that
rule. The result can be traced back to the state of the apparatus, to the reader, and to the
transport that fixed its scale.

\section{Bidirectional wave and temporal orientation}

A persistent route admits advanced and retarded channels. The reversal
of orientation acts on the word and the ledger; in the linear realization,
it exchanges the corresponding propagators when the Hamiltonian possesses the
declared symmetry. Bidirectional motion is expressed by a pair
\[
 (U_{\rm ret},U_{\rm adv})
\]
related by conjugation and time reversal.

The particle--antiparticle band realizes this duality as two orientations
of the same section. The dynamical relation can be expressed exactly.

\begin{proposition}[Conjugation of Orientation and Propagator Reversion]
Let \(C\) be an orthogonal involution on the quantum realification such that
\[
 C J_E C^{-1}=-J_E,
 \qquad
 C H C^{-1}=H.
\]
Then, for
\(U(t)=\exp(-J_EHt/\hbar_{\rm int})\),
\[
 C U(t)C^{-1}=U(-t).
\]
\end{proposition}

\begin{proof}
Withjugation commutes with the exponential series and transforms the generator
\(-J_EH\) in \(+J_EH\). The result is the exponential evaluated at
\(-t\).
\end{proof}

The word involution \(C_M=-\operatorname{rev}\) provides the genealogical
orientation datum. When its realization satisfies the two identities of the
proposition, the particle and antiparticle branches are represented by
propagators with conjugate temporal directions. The band geometry fixes the
orientation, and the Hamiltonian fixes the dynamics.

\section{Global expansion and self-consistency equation}

A genealogical cosmology can be formulated through an operator
\(\mathcal T\) that carries one boundary totality to the next:
\[
 B_{n+1}=\mathcal T(B_n).
\]
The locally induced response is
\[
 \mathcal I_{v,n}=\overline{\mathcal I}_v(B_n).
\]
A global self-inertial regime corresponds to a stable orbit or fixed point
of the transport operator, while expansion is recorded by a
scale cocycle \(a_{n+1}/a_n\). This formulation separates three questions:
existence of the dynamics, stability of the solution, and identification of the
cocycle with the observed cosmological factor.

\begin{resultbox}
Local inertia may be typed as the parallel transport of a section; its
Machian dependence is expressed as factorization through the global boundary
trace. The Ambivalence Principle links the areal and volumetric readings
of the same state. The observer enters through the coupling of frames that
makes measurement possible.
\end{resultbox}

\begin{statusbox}
The definition of self-inertia and the Machian factorization are recovered
results. The dimensional criterion and the measuring instrument follow from
the preceding chapters. The synthesis \eqref{eq:triptico-inercia-mach} is a
new formalization of a pre-existing authorial architecture. A quantitative
cosmological equation requires fixing \(\mathcal T\), its metric, and its boundary
data.
\end{statusbox}

The Machian attribution is characterized by the constancy of the local
response along the fibers of the global trace. Self-inertia is characterized
by parallel transport and by the fixed subspace of the holonomy. In a
cosmological realization, the operator \(\mathcal T\), its metric, and the scale
cocycle turn these structural conditions into a quantitative evolution of
the totality.
